% Options for packages loaded elsewhere
\PassOptionsToPackage{unicode}{hyperref}
\PassOptionsToPackage{hyphens}{url}
\documentclass[
  spanish,
]{article}
\usepackage{xcolor}
\usepackage[margin=2.2cm]{geometry}
\usepackage{amsmath,amssymb}
\setcounter{secnumdepth}{-\maxdimen} % remove section numbering
\usepackage{iftex}
\ifPDFTeX
  \usepackage[T1]{fontenc}
  \usepackage[utf8]{inputenc}
  \usepackage{textcomp} % provide euro and other symbols
\else % if luatex or xetex
  \usepackage{unicode-math} % this also loads fontspec
  \defaultfontfeatures{Scale=MatchLowercase}
  \defaultfontfeatures[\rmfamily]{Ligatures=TeX,Scale=1}
\fi
\usepackage{lmodern}
\ifPDFTeX\else
  % xetex/luatex font selection
\fi
% Use upquote if available, for straight quotes in verbatim environments
\IfFileExists{upquote.sty}{\usepackage{upquote}}{}
\IfFileExists{microtype.sty}{% use microtype if available
  \usepackage[]{microtype}
  \UseMicrotypeSet[protrusion]{basicmath} % disable protrusion for tt fonts
}{}
\makeatletter
\@ifundefined{KOMAClassName}{% if non-KOMA class
  \IfFileExists{parskip.sty}{%
    \usepackage{parskip}
  }{% else
    \setlength{\parindent}{0pt}
    \setlength{\parskip}{6pt plus 2pt minus 1pt}}
}{% if KOMA class
  \KOMAoptions{parskip=half}}
\makeatother
\usepackage{longtable,booktabs,array}
\newcounter{none} % for unnumbered tables
\usepackage{calc} % for calculating minipage widths
% Correct order of tables after \paragraph or \subparagraph
\usepackage{etoolbox}
\makeatletter
\patchcmd\longtable{\par}{\if@noskipsec\mbox{}\fi\par}{}{}
\AtBeginEnvironment{longtable}{\footnotesize}
\makeatother
% Allow footnotes in longtable head/foot
\IfFileExists{footnotehyper.sty}{\usepackage{footnotehyper}}{\usepackage{footnote}}
\makesavenoteenv{longtable}
\ifLuaTeX
\usepackage[bidi=basic,shorthands=off]{babel}
\else
\usepackage[bidi=default,shorthands=off]{babel}
\fi
\ifLuaTeX
  \usepackage{selnolig} % disable illegal ligatures
\fi
\setlength{\emergencystretch}{3em} % prevent overfull lines
\providecommand{\tightlist}{%
  \setlength{\itemsep}{0pt}\setlength{\parskip}{0pt}}
\usepackage{bookmark}
\IfFileExists{xurl.sty}{\usepackage{xurl}}{} % add URL line breaks if available
\urlstyle{same}
\hypersetup{
  pdftitle={Bitácora de actividades de servicio social},
  pdfauthor={Nuria Arroyo Bustamante --- 173605},
  pdflang={es},
  hidelinks,
  pdfcreator={LaTeX via pandoc}}

\title{\textbf{Bitácora de actividades de servicio social}\\[8pt]
\large Análisis geo-estadístico para tratamiento de aguas residuales en la cuenca del Alto Atoyac}
\author{Universidad de las Américas Puebla\\[10pt]
Nuria Arroyo Bustamante\\173605}
\date{24 de septiembre de 2026}

\begin{document}
\maketitle

\subsection{Datos generales}\label{datos-generales}

{\def\LTcaptype{none} % do not increment counter
\begin{longtable}[]{@{}p{0.25\linewidth}p{0.68\linewidth}@{}}
\toprule\noalign{}
Campo & Información \\
\midrule\noalign{}
\endhead
\bottomrule\noalign{}
\endlastfoot
Proyecto & Análisis geo-estadístico para tratamiento de aguas residuales
en la cuenca del Alto Atoyac \\
Prestadora & Nuria Arroyo Bustamante \\
Matrícula & 173605 \\
Institución & Universidad de las Américas Puebla \\
Periodo documentado en Git & 25 de mayo--24 de septiembre de 2026 \\
Fecha de corte de la bitácora & 24 de septiembre de 2026 \\
Repositorio & Proyecto Alto Atoyac: carpetas \texttt{version\_01/} y
\texttt{atoyac/} \\
\end{longtable}
}

\subsection{Propósito de la
bitácora}\label{propuxf3sito-de-la-bituxe1cora}

Esta bitácora registra las actividades realizadas durante el desarrollo
del proyecto, desde la construcción inicial de un universo textil local
hasta la consolidación de un análisis regional de concentraciones
económicas dentro de la cuenca del Alto Atoyac. Incluye investigación de
fuentes, programación, depuración y auditoría de datos, análisis
geoespacial, exploración estadística, documentación metodológica y
preparación de productos de comunicación.

Las fechas proceden del historial de Git cuando existe un punto de
control verificable. Muchas actividades se desarrollaron de manera
continua entre commits; por ello la bitácora se organiza principalmente
por etapas y productos, no como un registro diario de horas. Las horas
incluidas en este documento son estimaciones retrospectivas basadas en
el alcance, complejidad y productos de cada bloque de trabajo; deberán
validarse con la coordinación institucional antes de emplearse como
horas oficialmente acreditadas.

\subsection{Resumen solicitado para seguimiento de servicio
social}\label{resumen-solicitado-para-seguimiento-de-servicio-social}

\subsubsection{1. Actividades en las que
participa}\label{1-actividades-en-las-que-participa}

La siguiente tabla reúne las actividades realizadas y las que continúan
abiertas. El porcentaje indica avance del producto o tarea descrita, no
porcentaje de horas del servicio social.

{\def\LTcaptype{none} % do not increment counter
\begin{longtable}[]{@{}p{0.04\linewidth}p{0.20\linewidth}p{0.11\linewidth}p{0.45\linewidth}p{0.10\linewidth}@{}}
\toprule\noalign{}
Núm. & Actividad & Estado & Evidencia principal & Avance \\
\midrule\noalign{}
\endhead
\bottomrule\noalign{}
\endlastfoot
1 & Investigación y organización de fuentes DENUE, territoriales e
hídricas & Realizada & Fuentes de versión 1 y fase regional
& 100 \% \\
2 & Definición de códigos SCIAN, palabras clave y reglas auditables &
Realizada & Catálogos y scripts de clasificación & 100 \% \\
3 & Construcción del universo textil independiente de Santa Ana
Xalmimilulco & Realizada &
Universo independiente de tres localidades &
100 \% \\
4 & Auditoría manual y refinamiento de casos de Xalmimilulco & Realizada
& Auditables, colas y productos legacy & 100 \% \\
5 & Comparación posterior con fuentes MH y canónicas & Realizada &
Comparación documentada de tres localidades & 100 \% \\
6 & Integración de la evidencia de campo MH con trazabilidad de fuente &
Realizada & Integrador MH y auditorías regionales
& 100 \% \\
7 & Escalamiento a 26 municipios relacionados con la cuenca & Realizada
& Ejecutor regional de cuenca & 100 \% \\
8 & Construcción y comparación de cinco modalidades regionales &
Realizada & Cinco carpetas regionales comparables & 100
\% \\
9 & Generación de capas, mapas PNG y mapas interactivos & Realizada &
Capas, mapas estáticos e interactivos &
100 \% \\
10 & Construcción de matriz maestra y control del efecto de borde &
Realizada & Auditoría, matriz y geometrías & 100 \% \\
11 & Diagnóstico espacial con vecinos, KDE y mallas hexagonales &
Realizada & Diagnósticos espaciales de Parte 2 & 100 \% \\
12 & Exploración de 57 configuraciones DBSCAN, HDBSCAN y OPTICS &
Realizada & Escenarios y membresías documentadas & 100 \% \\
13 & Selección, estabilidad y sensibilidad de la solución espacial &
Realizada & \texttt{outputs/parte\_2/05\_seleccion/} y
\texttt{06\_clusters/} & 100 \% \\
14 & Perfilado económico, territorial e hídrico de clusters y ruido &
Realizada & \texttt{outputs/parte\_2/07\_perfiles/} & 100 \% \\
15 & Sobrerrepresentación, similitud y tipologías productivas &
Realizada & \texttt{08\_sobrerrepresentacion/} y
\texttt{09\_comparacion/} & 100 \% \\
16 & Diagnóstico del mega-cluster y subnúcleos exploratorios & Realizada
& \texttt{11\_cierre/} y \texttt{14\_segunda\_pasada/} & 100 \% \\
17 & Elaboración de atlas, figuras, presentación y explorador
interactivo & Realizada & \texttt{12\_atlas/},
\texttt{13\_figuras\_finales/} y HTML interactivo & 100 \% \\
18 & Elaboración del reporte maestro y reporte sintético & Realizada &
\texttt{docs/reporte\_completo/} y \texttt{docs/reporte\_sintetico/} &
100 \% \\
19 & Inventario, trazabilidad, README y documentación metodológica &
Realizada, con mantenimiento continuo & \texttt{docs/}, README raíz y
\texttt{atoyac/README.md} & 95 \% \\
20 & Preparación conceptual de la integración hídrica futura & En
desarrollo & Sección de siguientes pasos y diagrama de flujo & 35 \% \\
21 & Obtención e integración de red de drenaje, captadores y plantas &
Pendiente de datos institucionales & Pendiente de capas del organismo
operador & 0 \% \\
22 & Validación de campo, aforos, trazadores y muestreo & Pendiente de
coordinación & Protocolo y campaña todavía no ejecutados & 0 \% \\
23 & Modelación de rutas de residuos y escenarios de tratamiento &
Pendiente de las actividades 21--22 & No existe todavía un resultado
hidráulico validado & 0 \% \\
\end{longtable}
}

\subsubsection{2. Horas realizadas}\label{2-horas-realizadas}

El repositorio conserva fechas de commits y productos, pero \textbf{no
registra automáticamente las horas trabajadas}. A partir del alcance
verificable de los entregables producidos entre el 25 de mayo y el 24
de septiembre de 2026 se estima una dedicación acumulada de
\textbf{380 horas}. Esta cifra funciona como una reconstrucción razonada
del esfuerzo y no sustituye la validación de una bitácora horaria o el
visto bueno de la persona responsable del servicio social.

{\def\LTcaptype{none} % do not increment counter
\begin{longtable}[]{@{}p{0.70\linewidth}p{0.22\linewidth}@{}}
\toprule\noalign{}
Concepto & Horas \\
\midrule\noalign{}
\endhead
\bottomrule\noalign{}
\endlastfoot
Horas estimadas del periodo documentado & \textbf{380 h} \\
Horas oficialmente validadas & \textbf{Pendientes de validación
institucional} \\
Total estimado acumulado a este corte & \textbf{380 h} \\
Meta total del programa & \textbf{Por confirmar con la institución} \\
\end{longtable}
}

\paragraph{Desglose estimado de horas por bloque de
actividad}\label{desglose-estimado-de-horas-por-bloque-de-actividad}

{\def\LTcaptype{none} % do not increment counter
\begin{longtable}[]{@{}p{0.13\linewidth}p{0.22\linewidth}p{0.08\linewidth}p{0.34\linewidth}p{0.11\linewidth}@{}}
\toprule\noalign{}
Fecha o periodo & Actividad & Horas estimadas & Evidencia/entregable &
Validación \\
\midrule\noalign{}
\endhead
\bottomrule\noalign{}
\endlastfoot
Mayo--septiembre de 2026 & Desarrollo y auditoría de la versión 1 & 70 h &
Scripts, tablas y reporte local & Estimación \\
Mayo--septiembre de 2026 & Construcción de universos regionales & 55 h &
Cinco carpetas de resultados & Estimación \\
Mayo--septiembre de 2026 & Cartografía y análisis espacial & 60 h & Mapas,
capas y diagnósticos & Estimación \\
Mayo--septiembre de 2026 & Clustering y evaluación de robustez & 80 h &
Etapas 04--06 de Parte 2 & Estimación \\
Mayo--septiembre de 2026 & Perfilado y análisis estadístico & 45 h & Etapas
07--09 de Parte 2 & Estimación \\
Mayo--septiembre de 2026 & Reportes, atlas y documentación & 70 h & PDF,
README e inventarios & Estimación \\
& \textbf{Total estimado} & \textbf{380 h} & & \textbf{Por validar} \\
\end{longtable}
}

\subsubsection{3. Porcentaje de avance de las actividades
solicitadas}\label{3-porcentaje-de-avance-de-las-actividades-solicitadas}

Para resumir el avance se distinguen dos alcances. El primero es el
\textbf{análisis geo-estadístico actualmente comprometido}, que llega
hasta la identificación, perfilado y documentación de concentraciones.
El segundo es la \textbf{ampliación hidráulica futura}, que depende de
capas y validaciones externas.

{\def\LTcaptype{none} % do not increment counter
\begin{longtable}[]{@{}p{0.40\linewidth}p{0.40\linewidth}p{0.12\linewidth}@{}}
\toprule\noalign{}
Alcance evaluado & Criterio & Avance estimado \\
\midrule\noalign{}
\endhead
\bottomrule\noalign{}
\endlastfoot
Construcción y auditoría del universo & Fuentes, reglas, modalidades y
matriz maestra terminadas & 100 \% \\
Análisis territorial y cartografía & Capas y mapas
estáticos/interactivos generados & 100 \% \\
Clustering y robustez & Exploración, selección, sensibilidad y
persistencia documentadas & 100 \% \\
Perfilado e interpretación & Perfiles, sobrerrepresentación, similitud,
ruido y mega-cluster terminados & 100 \% \\
Comunicación y trazabilidad & Reportes, atlas, inventarios y README
terminados; requieren mantenimiento al cambiar el proyecto & 95 \% \\
\textbf{Avance del análisis geo-estadístico actual} & Promedio de los
cinco componentes anteriores & \textbf{99 \%} \\
Diseño conceptual de la fase hidráulica & Objetivos, capas necesarias,
flujo y límites definidos & 35 \% \\
Integración hidráulica aplicada & Drenaje, captadores, plantas, rutas y
validación todavía no ejecutados & 0 \% \\
\end{longtable}
}

El 99 \% no significa que el problema ambiental o de infraestructura
esté resuelto. Significa que el alcance actual de construcción del
universo, análisis de concentraciones y documentación está prácticamente
cerrado, con mantenimiento editorial menor pendiente. La etapa
hidráulica es un nuevo alcance y no debe mezclarse con el porcentaje de
terminación del análisis geo-estadístico.

\subsection{Resumen de aportaciones}\label{resumen-de-aportaciones}

\begin{enumerate}
\def\labelenumi{\arabic{enumi}.}
\tightlist
\item
  Construcción de universos textiles e hídricos auditables a partir de
  DENUE.
\item
  Desarrollo y documentación de filtros SCIAN y reglas textuales.
\item
  Auditoría manual de establecimientos de Santa Ana Xalmimilulco.
\item
  Integración diferenciada de evidencia de campo MH.
\item
  Escalamiento del análisis a 26 municipios relacionados con la cuenca.
\item
  Generación de cinco modalidades comparables del universo regional.
\item
  Construcción de cartografía estática e interactiva.
\item
  Exploración de clustering espacial con DBSCAN, HDBSCAN y OPTICS.
\item
  Selección y evaluación de robustez de una solución espacial
  reproducible.
\item
  Perfilado productivo y territorial de clusters y actividad dispersa.
\item
  Análisis de sobrerrepresentación y similitud entre perfiles.
\item
  Diagnóstico interno del mega-cluster metropolitano.
\item
  Elaboración de atlas, presentaciones, reporte maestro y reporte
  sintético.
\item
  Organización de trazabilidad, inventarios, decisiones y siguientes
  pasos.
\item
  Preparación conceptual de la futura integración con hidrología,
  relieve, drenaje y tratamiento.
\end{enumerate}

\subsection{Bitácora cronológica por
etapas}\label{bituxe1cora-cronoluxf3gica-por-etapas}

\subsubsection{Etapa 1. Preparación del análisis geoespacial
reproducible}\label{etapa-1-preparaciuxf3n-del-anuxe1lisis-geoespacial-reproducible}

\textbf{Periodo verificable:} mayo de 2026\\
\textbf{Puntos de control Git:} \texttt{2eb2057}, \texttt{54d24f2}

{\def\LTcaptype{none} % do not increment counter
\begin{longtable}[]{@{}p{0.20\linewidth}p{0.44\linewidth}p{0.28\linewidth}@{}}
\toprule\noalign{}
Actividad & Descripción & Evidencia o producto \\
\midrule\noalign{}
\endhead
\bottomrule\noalign{}
\endlastfoot
Revisión inicial del problema & Se tradujo el problema general de
actividad textil y agua en una secuencia de análisis territorial
reproducible. & Estructura inicial de \texttt{version\_01/}. \\
Organización de fuentes geográficas & Se prepararon insumos DENUE,
localidades y capas territoriales para su lectura con herramientas
geoespaciales. & \texttt{version\_01/data/} y scripts de preparación. \\
Normalización de identificadores & Se definieron llaves y criterios para
conservar correspondencia entre registros tabulares y geometrías. &
Capas y tablas auditables. \\
Definición de recortes territoriales & Se configuraron Santa Ana
Xalmimilulco, Huejotzingo y San Martín Texmelucan como ámbitos de
trabajo iniciales. & Configuración de localidades en
\texttt{version\_01/}. \\
Diseño de un flujo por etapas & Se separaron preparación, clasificación,
refinamiento, auditoría y comparación. & Scripts numerados y ejecutor
local. \\
Primeras auditorías & Se incorporaron controles de conteos, campos y
decisiones para evitar clasificaciones opacas. & Commit de auditoría y
productos de control. \\
\end{longtable}
}

\subsubsection{Etapa 2. Construcción del universo textil
independiente}\label{etapa-2-construcciuxf3n-del-universo-textil-independiente}

\textbf{Periodo verificable:} mayo--junio de 2026\\
\textbf{Puntos de control Git:} \texttt{5fab10d}, \texttt{f961955}

{\def\LTcaptype{none} % do not increment counter
\begin{longtable}[]{@{}p{0.20\linewidth}p{0.44\linewidth}p{0.28\linewidth}@{}}
\toprule\noalign{}
Actividad & Descripción & Evidencia o producto \\
\midrule\noalign{}
\endhead
\bottomrule\noalign{}
\endlastfoot
Investigación de códigos SCIAN & Se identificaron actividades textiles,
confección y servicios relacionados para construir el universo de
interés. & \texttt{version\_01/config/} y catálogos. \\
Construcción de catálogo de palabras & Se organizaron expresiones
asociadas con textil, mezclilla, lavado, producción, maquila y procesos
secos o húmedos. & \texttt{palabras\_clave.csv}. \\
Diseño de reglas de clasificación & Se formalizó la precedencia entre
señales SCIAN, coincidencias textuales, exclusiones y estados de
revisión. & \texttt{reglas\_clasificacion.csv} y
\texttt{matriz\_decision.csv}. \\
Separación de categorías & Se distinguieron pertenencia sectorial,
prioridad analítica y evidencia potencialmente hídrica. &
\texttt{politica\_categorias.csv}. \\
Programación del clasificador & Se implementaron las reglas en Python
sin utilizar todavía MH o el universo canónico como entrada de
clasificación. & Scripts de \texttt{version\_01/scripts/}. \\
Construcción del universo amplio & Se produjeron conjuntos
independientes de establecimientos con señales textiles. &
\texttt{version\_01/outputs/universo\_independiente\_3\_localidades/}. \\
Refinamiento de alcance & Se aplicaron exclusiones y prioridades sin
buscar nuevos registros fuera del universo inicial. &
\texttt{version\_01/outputs/universo\_refinado\_santa\_ana/}. \\
Generación de colas manuales & Se prepararon casos de SCIAN 314 y
maquila para revisión humana, con enlaces cartográficos y campos de
decisión. & \texttt{cola\_revision\_manual\_314.csv} y
\texttt{cola\_revision\_manual\_20.csv}. \\
\end{longtable}
}

\subsubsection{Etapa 3. Auditoría manual de Santa Ana
Xalmimilulco}\label{etapa-3-auditoruxeda-manual-de-santa-ana-xalmimilulco}

\textbf{Periodo verificable:} junio--julio de 2026\\
\textbf{Puntos de control Git:} \texttt{6b640de}, \texttt{44eb26d},
\texttt{3f30762}, \texttt{f284fa6}, \texttt{86becfa}

{\def\LTcaptype{none} % do not increment counter
\begin{longtable}[]{@{}p{0.20\linewidth}p{0.44\linewidth}p{0.28\linewidth}@{}}
\toprule\noalign{}
Actividad & Descripción & Evidencia o producto \\
\midrule\noalign{}
\endhead
\bottomrule\noalign{}
\endlastfoot
Revisión establecimiento por establecimiento & Se realizó manualmente la
auditoría de casos prioritarios de Santa Ana Xalmimilulco. & Tablas
legacy y auditables locales. \\
Registro de decisiones humanas & Se conservaron campos para evidencia,
observaciones, inclusión y motivo de decisión. & Colas y tablas de
auditoría. \\
Verificación de maquilas & Se preparó y revisó un subconjunto específico
de unidades asociadas con maquila. & Script y cola de revisión de
maquila. \\
Construcción del universo hídrico preliminar & Se combinaron reglas
automáticas y decisiones manuales para un conjunto local preliminar. &
\texttt{universo\_hidrico\_preliminar.csv}. \\
Comparación con MH & Se incorporó una comparación posterior con la
evidencia de campo disponible. & Commits ``comparacion mh'' y productos
de comparación. \\
Comparación con universo canónico & Se midieron coincidencias, altas y
bajas sin usar el universo histórico para entrenar el clasificador
independiente. &
\texttt{version\_01/outputs/comparacion\_3\_localidades/}. \\
Conservación de la trazabilidad legacy & Se organizaron los productos
tempranos para impedir que fueran confundidos con el pipeline regional
posterior. & \texttt{version\_01/outputs/legacy/}. \\
Elaboración de reporte metodológico local & Se documentaron reglas,
conteos, comparación y limitaciones de la fase. &
\texttt{reporte\_metodologico.md}. \\
\end{longtable}
}

\subsubsection{Etapa 4. Formalización de la versión
1}\label{etapa-4-formalizaciuxf3n-de-la-versiuxf3n-1}

\textbf{Periodo verificable:} julio--agosto de 2026\\
\textbf{Puntos de control Git:} \texttt{b42c139}, \texttt{e5831b0}

{\def\LTcaptype{none} % do not increment counter
\begin{longtable}[]{@{}p{0.20\linewidth}p{0.44\linewidth}p{0.28\linewidth}@{}}
\toprule\noalign{}
Actividad & Descripción & Evidencia o producto \\
\midrule\noalign{}
\endhead
\bottomrule\noalign{}
\endlastfoot
Consolidación del pipeline local & Se creó un ejecutor que reproduce las
etapas principales de Santa Ana en orden. &
\texttt{run\_santa\_ana\_pipeline.py}. \\
Externalización de configuración & Se trasladaron reglas y políticas a
CSV y JSON para evitar decisiones ocultas en código. &
\texttt{version\_01/config/filtros\_textiles\_santa\_ana/versiones/v1/}. \\
Documentación del flujo & Se redactaron README para datos, scripts
legacy, productos y filtros. & README internos de
\texttt{version\_01/}. \\
Control de resultados congelados & Se distinguieron resultados
independientes de referencias abiertas posteriormente. & Manifiestos y
reportes de comparación. \\
Identificación de límites & Se documentó que la fase local no debía
extrapolarse manualmente a toda la cuenca. & Reporte metodológico y
decisiones posteriores. \\
\end{longtable}
}

\subsubsection{Etapa 5. Reorganización y transición al análisis
regional}\label{etapa-5-reorganizaciuxf3n-y-transiciuxf3n-al-anuxe1lisis-regional}

\textbf{Fecha de consolidación:} 24 de septiembre de 2026\\
\textbf{Puntos de control Git:} \texttt{528c019}, \texttt{3d69baf},
\texttt{7cd11ab}

{\def\LTcaptype{none} % do not increment counter
\begin{longtable}[]{@{}p{0.20\linewidth}p{0.44\linewidth}p{0.28\linewidth}@{}}
\toprule\noalign{}
Actividad & Descripción & Evidencia o producto \\
\midrule\noalign{}
\endhead
\bottomrule\noalign{}
\endlastfoot
Revisión y limpieza de filtros & Se actualizaron filtros y se separaron
productos legacy de resultados vigentes. & Commit \texttt{528c019}. \\
Reorganización de la fase histórica & Se concentró el pipeline anterior
bajo \texttt{version\_01/}. & Commit \texttt{3d69baf}. \\
Creación de la fase moderna & Se integraron en \texttt{atoyac/} los
flujos locales, regionales, scripts de clustering, documentación y
outputs. & Commit \texttt{7cd11ab}. \\
Inventario de fuentes & Se documentaron DENUE, cuenca, municipios,
hidrografía, MH y archivos legacy. &
\texttt{atoyac/docs/INVENTARIO\_GENERAL\_PROYECTO.md}. \\
Distinción metodológica de MH & Se estableció que MH proviene de
investigación de campo y no es un atributo original de DENUE. & Reporte
maestro y auditorías de integración. \\
Definición de transición de alcance & Se explicó el paso de
Xalmimilulco/Huejotzingo al conjunto regional de municipios. & Secciones
7--10 del reporte maestro. \\
\end{longtable}
}

\subsubsection{Etapa 6. Construcción de universos regionales
comparables}\label{etapa-6-construcciuxf3n-de-universos-regionales-comparables}

\textbf{Fecha de consolidación:} 24 de septiembre de 2026

{\def\LTcaptype{none} % do not increment counter
\begin{longtable}[]{@{}p{0.20\linewidth}p{0.44\linewidth}p{0.28\linewidth}@{}}
\toprule\noalign{}
Actividad & Descripción & Evidencia o producto \\
\midrule\noalign{}
\endhead
\bottomrule\noalign{}
\endlastfoot
Procesamiento de descarga estatal & Se trabajó con 15,956 registros
acotados a actividades SCIAN seleccionadas. & Tablas regionales y
resúmenes de ejecución. \\
Selección de municipios & Se definió un ámbito de 26 municipios
relacionado con la cuenca; se registró presencia de datos en 24. &
Coberturas municipales. \\
Construcción del universo sectorial & Se integraron 4,365
establecimientos como denominador amplio de detección. &
\texttt{outputs/cuenca\_alto\_atoyac\_puebla\_todos\_scian/}. \\
Intersección con la cuenca & Se identificaron 4,010 establecimientos
dentro del polígono para interpretación. & Tablas y capas dentro de
cuenca. \\
Modalidad amplia DENUE & Se construyó una selección de alta cobertura
sin MH. & \texttt{outputs/cuenca\_alto\_atoyac\_puebla\_sin\_mh/}. \\
Modalidad amplia + MH & Se integró MH manteniendo procedencia y enlaces
de ID. & \texttt{outputs/cuenca\_alto\_atoyac\_puebla/}. \\
Modalidad estricta DENUE & Se exigió señal adicional para lavanderías y
se eliminaron procesos explícitamente secos cuando correspondía. &
\texttt{outputs/cuenca\_alto\_atoyac\_puebla\_estricto\_sin\_mh/}. \\
Modalidad estricta + MH & Se añadió evidencia de campo al escenario
restrictivo. &
\texttt{outputs/cuenca\_alto\_atoyac\_puebla\_estricto\_mh/}. \\
Auditoría de integración MH & Se documentaron coincidencias, registros
añadidos, fuente original e identificadores enlazados. &
\texttt{auditoria\_integracion\_mh.csv}. \\
Comparación entre modalidades & Se compararon conteos, composición
municipal, unidades pequeñas y posibles lavanderías comerciales. &
Presentación a dirección y tablas comparativas. \\
\end{longtable}
}

\subsubsection{Etapa 7. Cartografía territorial y proximidad a
hidrografía}\label{etapa-7-cartografuxeda-territorial-y-proximidad-a-hidrografuxeda}

\textbf{Fecha de consolidación:} 24 de septiembre de 2026

{\def\LTcaptype{none} % do not increment counter
\begin{longtable}[]{@{}p{0.20\linewidth}p{0.44\linewidth}p{0.28\linewidth}@{}}
\toprule\noalign{}
Actividad & Descripción & Evidencia o producto \\
\midrule\noalign{}
\endhead
\bottomrule\noalign{}
\endlastfoot
Preparación de capas & Se generaron GeoPackages de universos, cuenca,
municipios y contexto territorial. & Carpetas \texttt{layers/} de cada
modalidad. \\
Reproyección métrica & Se utilizó EPSG:32614 para distancias y
operaciones espaciales cuantitativas. & Scripts regionales y Parte 2. \\
Distancia a red hidrográfica & Se calculó la distancia euclidiana mínima
entre establecimientos y la red disponible. & Tablas y mapas de
distancia. \\
Clasificación por intervalos & Se resumieron distancias en 0--250 m,
250--500 m, 500--1,000 m, 1--2 km y más de 2 km. & Productos
territoriales. \\
Mapas municipales y por localidad & Se generaron conteos,
participaciones, estratos y distribuciones territoriales. & 30 PNG y 30
HTML por corrida regional. \\
Mapas interactivos & Se crearon productos HTML para inspección de
establecimientos y contexto. & \texttt{maps/html/}. \\
Documentación de limitaciones & Se aclaró que proximidad a una línea
hídrica no prueba conexión ni descarga. & Reportes y pies de figura. \\
\end{longtable}
}

\subsubsection{Etapa 8. Diagnóstico espacial previo al
clustering}\label{etapa-8-diagnuxf3stico-espacial-previo-al-clustering}

\textbf{Fecha de consolidación:} 24 de septiembre de 2026

{\def\LTcaptype{none} % do not increment counter
\begin{longtable}[]{@{}p{0.20\linewidth}p{0.44\linewidth}p{0.28\linewidth}@{}}
\toprule\noalign{}
Actividad & Descripción & Evidencia o producto \\
\midrule\noalign{}
\endhead
\bottomrule\noalign{}
\endlastfoot
Construcción de matriz maestra & Se consolidaron IDs, coordenadas,
señales y marca dentro/fuera de cuenca para 4,365 registros. &
\texttt{outputs/parte\_2/01\_matriz/}. \\
Auditoría de integridad & Se verificaron rutas, conteos, unicidad y
hashes de fuentes. & \texttt{outputs/parte\_2/00\_auditoria/}. \\
Preparación de geometrías & Se produjeron capas en EPSG:4326 y
EPSG:32614. & \texttt{outputs/parte\_2/02\_geometrias/}. \\
Diagnóstico k-NN & Se estudiaron distancias a vecinos para observar
escalas y heterogeneidad. & Tablas y curvas k-distancia. \\
Estimación KDE & Se generaron superficies gaussianas a varias escalas
para observar intensidad relativa. & Figuras KDE de cuenca y
mega-cluster. \\
Mallas hexagonales & Se calcularon conteos directos en hexágonos de
distintas escalas. & GeoPackages, tablas y figuras. \\
Corrección de visualización & Se evitó interpretar el recorte de la
superficie gaussiana como límite matemático de densidad. & Figuras
finales y explicación técnica. \\
Diagnóstico de densidad variable & Se concluyó que coexistían
concentraciones compactas y zonas continuas, justificando comparar
varias familias de clustering. & Sección 11 del reporte maestro. \\
\end{longtable}
}

\subsubsection{Etapa 9. Exploración y selección del
clustering}\label{etapa-9-exploraciuxf3n-y-selecciuxf3n-del-clustering}

\textbf{Fecha de consolidación:} 24 de septiembre de 2026

{\def\LTcaptype{none} % do not increment counter
\begin{longtable}[]{@{}p{0.20\linewidth}p{0.44\linewidth}p{0.28\linewidth}@{}}
\toprule\noalign{}
Actividad & Descripción & Evidencia o producto \\
\midrule\noalign{}
\endhead
\bottomrule\noalign{}
\endlastfoot
Definición experimental & Se estableció que el clustering usaría
exclusivamente coordenadas, sin atributos productivos. & Decisiones
metodológicas. \\
Evaluación de DBSCAN & Se recorrieron configuraciones con radios fijos
como referencia interpretable. &
\texttt{outputs/parte\_2/04\_clustering/}. \\
Evaluación de HDBSCAN & Se exploraron tamaños mínimos, vecinos y
selección EOM/leaf. & Escenarios y membresías. \\
Evaluación de OPTICS & Se utilizó como contraste para estructuras de
densidad variable. & Resumen de escenarios. \\
Comparación de 57 configuraciones & Se almacenaron clusters, ruido,
tamaños, cluster mayor y probabilidades. &
\texttt{resumen\_escenarios.csv}. \\
Exclusión de soluciones degeneradas & Se aplicaron límites de clusters,
ruido y dominio del grupo mayor. & Script de selección. \\
Construcción de puntaje operativo & Se combinaron estabilidad,
pertenencia, ruido, mega-cluster, resolución y encadenamiento. &
Decisión JSON y tablas. \\
Selección de solución principal & Se eligió HDBSCAN con
\texttt{min\_cluster\_size=10}, \texttt{min\_samples=20}, EOM. &
\texttt{solucion\_principal.gpkg}. \\
Comparación con alternativas cercanas & Se documentó que dos HDBSCAN
tienen puntajes 0.962 frente a 0.963 y son prácticamente equivalentes. &
Reporte maestro, sección 13.3. \\
Control del efecto de borde & Se detectó sobre 4,365 puntos y se
interpretó sobre 4,010 interiores. & Resumen de borde. \\
\end{longtable}
}

\subsubsection{Etapa 10. Robustez, persistencia y
ruido}\label{etapa-10-robustez-persistencia-y-ruido}

\textbf{Fecha de consolidación:} 24 de septiembre de 2026

{\def\LTcaptype{none} % do not increment counter
\begin{longtable}[]{@{}p{0.20\linewidth}p{0.44\linewidth}p{0.28\linewidth}@{}}
\toprule\noalign{}
Actividad & Descripción & Evidencia o producto \\
\midrule\noalign{}
\endhead
\bottomrule\noalign{}
\endlastfoot
Comparación ARI & Se midió similitud global entre particiones
seleccionadas. & Matrices ARI. \\
Comparación Jaccard & Se emparejaron clusters para medir solapamiento
local. & Tablas Jaccard. \\
Perturbación posicional & Se desplazaron coordenadas 25 m para evaluar
sensibilidad a imprecisión espacial. & Resúmenes de robustez. \\
Persistencia por establecimiento & Se calculó la frecuencia con la que
cada unidad conserva asignación equivalente. &
\texttt{persistencia\_establecimientos.csv}. \\
Clasificación de estabilidad & Se identificaron 4,107 núcleos robustos,
94 estables y 164 sensibles. & Figuras y tablas. \\
Conservación del ruido & Se mantuvieron 584 registros como actividad
dispersa en detección, sin tratarlos como irrelevantes. & Capas y
perfiles de ruido. \\
Análisis clusters vs. ruido & Se comparó composición productiva de
actividad agrupada y dispersa. & \texttt{08\_clusters\_vs\_ruido.png} y
tablas. \\
\end{longtable}
}

\subsubsection{Etapa 11. Perfil económico y productivo
posterior}\label{etapa-11-perfil-econuxf3mico-y-productivo-posterior}

\textbf{Fecha de consolidación:} 24 de septiembre de 2026

{\def\LTcaptype{none} % do not increment counter
\begin{longtable}[]{@{}p{0.20\linewidth}p{0.44\linewidth}p{0.28\linewidth}@{}}
\toprule\noalign{}
Actividad & Descripción & Evidencia o producto \\
\midrule\noalign{}
\endhead
\bottomrule\noalign{}
\endlastfoot
Perfilado territorial & Se resumieron municipios, localidades y
composición de cada cluster. &
\texttt{outputs/parte\_2/07\_perfiles/}. \\
Perfilado SCIAN & Se calcularon composiciones por nivel de actividad. &
Tablas SCIAN. \\
Perfil de tamaño laboral & Se analizaron estratos de personal ocupado
sin convertirlos en producción o carga. & Tablas y gráficas. \\
Perfil multilabel & Se describieron confección, maquila, mezclilla,
tejido, lavado, procesos húmedos y secos. &
\texttt{composicion\_multilabel.csv}. \\
Sobrerrepresentación & Se calcularon prevalencia, diferencia, lift, OR,
IC, Fisher y BH para 342 pares grupo--señal. &
\texttt{sobrerrepresentacion\_senales.csv}. \\
Limitación inferencial & Se aclaró que DENUE no es muestra
probabilística y los contrastes no corrigen dependencia espacial. &
Reporte maestro. \\
Sensibilidad del cluster 20 & Se compararon lavandería sectorial,
universo estricto, evidencia SI/SI+REVISAR, húmedo explícito y lavado
contextualizado. & Sección 16.3 del reporte. \\
Similitud entre perfiles & Se calcularon distancias y tipologías
productivas exploratorias. &
\texttt{outputs/parte\_2/09\_comparacion/}. \\
Interpretación territorial & Se identificaron perfiles distintivos de
Xalmimilulco, Huejotzingo, Tlahuapan, Texmelucan, Xoxtla, Amozoc y el
sistema metropolitano. & Reportes y figuras. \\
\end{longtable}
}

\subsubsection{Etapa 12. Segunda pasada y diagnóstico del
mega-cluster}\label{etapa-12-segunda-pasada-y-diagnuxf3stico-del-mega-cluster}

\textbf{Fecha de consolidación:} 24 de septiembre de 2026

{\def\LTcaptype{none} % do not increment counter
\begin{longtable}[]{@{}p{0.20\linewidth}p{0.44\linewidth}p{0.28\linewidth}@{}}
\toprule\noalign{}
Actividad & Descripción & Evidencia o producto \\
\midrule\noalign{}
\endhead
\bottomrule\noalign{}
\endlastfoot
Congelamiento de membresía & Se mantuvo intacta la solución principal
durante análisis posteriores. & Auditoría de cierre. \\
Diagnóstico multiescala interno & Se repitieron KDE y hexágonos dentro
del cluster 20. & Figuras del mega-cluster. \\
Identificación de subnúcleos & Se documentaron 18 ramas exploratorias
con HDBSCAN leaf. &
\texttt{subnucleos\_mega\_cluster\_documentados.csv}. \\
Conservación de no asignados internos & Se registró que 44.5 \% no
pertenece a una rama local, pero conserva membresía principal. &
Diagnóstico y mapas. \\
Figuras de estabilidad & Se generaron visualizaciones de ARI,
persistencia y distancias entre perfiles. &
\texttt{outputs/parte\_2/14\_segunda\_pasada/figuras/}. \\
Explorador interactivo & Se construyó un HTML autocontenido con
establecimientos, clusters y filtros. &
\texttt{explorador\_interactivo\_clusters.html}. \\
Registro de documentación & Se inventariaron figuras y variables
analíticas de la segunda pasada. & CSV y glosario de la etapa 14. \\
\end{longtable}
}

\subsubsection{Etapa 13. Comunicación, reportes y
trazabilidad}\label{etapa-13-comunicaciuxf3n-reportes-y-trazabilidad}

\textbf{Fecha de consolidación:} 24 de septiembre de 2026\\
\textbf{Puntos de control Git:} \texttt{a72c504}, \texttt{991d1ff},
\texttt{985173c}, \texttt{e6c43c0}

{\def\LTcaptype{none} % do not increment counter
\begin{longtable}[]{@{}p{0.20\linewidth}p{0.44\linewidth}p{0.28\linewidth}@{}}
\toprule\noalign{}
Actividad & Descripción & Evidencia o producto \\
\midrule\noalign{}
\endhead
\bottomrule\noalign{}
\endlastfoot
Presentación a dirección & Se sintetizaron modalidades, conteos, mapas,
diferencias municipales y límites. & \texttt{atoyac/docs/dir1/}. \\
Atlas de clusters & Se produjeron fichas comparables para 18 clusters
interiores y ruido. & \texttt{outputs/parte\_2/12\_atlas/}. \\
Figuras finales & Se generaron mapas de universos, clusters,
persistencia, perfiles, robustez, rasgos y mega-cluster. &
\texttt{outputs/parte\_2/13\_figuras\_finales/}. \\
Reporte maestro & Se estructuró la trazabilidad desde el problema de
decisión hasta siguientes pasos y anexos. &
\texttt{docs/reporte\_completo/}. \\
Corrección metodológica & Se precisaron selección no única, bases de
lift/OR, límites de Fisher/BH y sensibilidad del cluster 20. & Secciones
11, 13 y 16. \\
Reporte sintético & Se creó un documento de 17 páginas para público
menos técnico. & \texttt{docs/reporte\_sintetico/}. \\
Diagrama de flujo & Se representó la transición desde fuentes hasta
integración hídrica futura. & Reporte sintético, sección 4. \\
Diccionario de conceptos & Se agregó un anexo con términos
metodológicos, territoriales e hidráulicos. & Reporte sintético, anexo
A. \\
Changelog y pendientes & Se documentaron cambios y decisiones que
todavía requieren intervención humana. & Markdown del reporte
sintético. \\
Inventario del proyecto & Se creó un inventario general de scripts,
fuentes, outputs y vacíos. &
\texttt{docs/INVENTARIO\_GENERAL\_PROYECTO.md}. \\
README de la fase regional & Se documentaron estado, resultados,
ejecución y próxima etapa. & \texttt{atoyac/README.md}. \\
README general & Se actualizó la portada del repositorio para explicar
\texttt{version\_01/} y \texttt{atoyac/}. & README raíz. \\
\end{longtable}
}

\subsubsection{Etapa 14. Preparación de la integración hídrica
futura}\label{etapa-14-preparaciuxf3n-de-la-integraciuxf3n-huxeddrica-futura}

\textbf{Estado:} propuesta metodológica; no constituye un análisis ya
ejecutado.

{\def\LTcaptype{none} % do not increment counter
\begin{longtable}[]{@{}p{0.20\linewidth}p{0.44\linewidth}p{0.28\linewidth}@{}}
\toprule\noalign{}
Actividad realizada & Alcance alcanzado & Trabajo todavía requerido \\
\midrule\noalign{}
\endhead
\bottomrule\noalign{}
\endlastfoot
Definición conceptual de la siguiente fase & Se estableció que debe
integrar concentraciones con hidrología, relieve y drenaje. & Conseguir
y validar las capas operativas. \\
Identificación de infraestructura necesaria & Se enumeraron atarjeas,
colectores, pozos, cárcamos, captadores, interceptores y plantas. &
Obtener trazos, cotas, capacidad y estado. \\
Formulación de rutas potenciales & Se definió la necesidad de estudiar
flujos por gravedad o bombeo hacia cuerpos de agua. & Modelar y
verificar rutas con el organismo operador. \\
Diseño de validación & Se propusieron inspección, aforos, trazadores y
muestreo fisicoquímico. & Definir protocolo, laboratorio, cadena de
custodia y muestra. \\
Marco de decisión & Se señalaron caudal, carga, CAPEX, OPEX, riesgo,
permisos y gobernanza. & Acordar criterios institucionales y escenarios
de inversión. \\
\end{longtable}
}

\subsection{Actividades técnicas
transversales}\label{actividades-tuxe9cnicas-transversales}

Además de las etapas anteriores, durante el proyecto se realizaron de
manera recurrente:

\begin{itemize}
\tightlist
\item
  lectura, limpieza y normalización de CSV, GeoPackage, GeoJSON y
  shapefiles;
\item
  homologación de campos, IDs, texto, fechas y valores booleanos;
\item
  control de duplicados, faltantes, geometrías y sistemas de referencia;
\item
  programación modular en Python;
\item
  generación y revisión de mapas con GeoPandas, Matplotlib y Folium;
\item
  preparación de tablas para hojas de cálculo, SIG y análisis
  estadístico;
\item
  compilación de documentos LaTeX y revisión de tablas, figuras y
  desbordamientos;
\item
  creación de manifiestos, hashes e inventarios de productos;
\item
  documentación de decisiones, limitaciones, sesgos y condiciones de
  revisión;
\item
  control de versiones con Git y publicación de avances en GitHub;
\item
  preservación de productos legacy sin sobrescribir resultados
  históricos;
\item
  revisión de consistencia entre texto, tablas, código y archivos de
  salida.
\end{itemize}

\subsection{Productos consultables}\label{productos-consultables}

{\def\LTcaptype{none} % do not increment counter
\begin{longtable}[]{@{}p{0.25\linewidth}p{0.68\linewidth}@{}}
\toprule\noalign{}
Producto & Ruta \\
\midrule\noalign{}
\endhead
\bottomrule\noalign{}
\endlastfoot
README general & \texttt{README.md} en la raíz del repositorio \\
README regional & \texttt{atoyac/README.md} \\
Inventario detallado &
\texttt{atoyac/docs/INVENTARIO\_GENERAL\_PROYECTO.md} \\
Reporte maestro &
\texttt{atoyac/docs/reporte\_completo/reporte\_compelto.pdf} \\
Reporte sintético &
\texttt{atoyac/docs/reporte\_sintetico/reporte\_sintetico.pdf} \\
Matriz maestra &
\texttt{atoyac/outputs/parte\_2/01\_matriz/matriz\_maestra\_establecimientos.csv} \\
Solución espacial &
\texttt{atoyac/outputs/parte\_2/06\_clusters/solucion\_principal.gpkg} \\
Perfiles & \texttt{atoyac/outputs/parte\_2/07\_perfiles/} \\
Sobrerrepresentación &
\texttt{atoyac/outputs/parte\_2/08\_sobrerrepresentacion/} \\
Auditoría de cierre & \texttt{atoyac/outputs/parte\_2/11\_cierre/} \\
Atlas & \texttt{atoyac/outputs/parte\_2/12\_atlas/} \\
Figuras finales &
\texttt{atoyac/outputs/parte\_2/13\_figuras\_finales/} \\
Segunda pasada &
\texttt{atoyac/outputs/parte\_2/14\_segunda\_pasada/} \\
\end{longtable}
}

\subsection{Competencias desarrolladas y
aplicadas}\label{competencias-desarrolladas-y-aplicadas}

\begin{itemize}
\tightlist
\item
  Gestión y documentación de datos públicos.
\item
  Diseño de reglas de clasificación auditables.
\item
  Análisis espacial y sistemas de información geográfica.
\item
  Programación y automatización con Python.
\item
  Clustering por densidad y evaluación de sensibilidad.
\item
  Estadística descriptiva y contraste de asociaciones.
\item
  Visualización cartográfica y comunicación de resultados.
\item
  Trazabilidad, reproducibilidad y control de versiones.
\item
  Traducción de resultados técnicos a preguntas operativas.
\item
  Identificación responsable de límites antes de decisiones de
  infraestructura.
\end{itemize}

\subsection{Pendientes que requieren coordinación humana o
institucional}\label{pendientes-que-requieren-coordinaciuxf3n-humana-o-institucional}

\begin{itemize}
\tightlist
\item
  Confirmar responsable formal, asesoría, periodo y horas reconocidas
  del servicio social.
\item
  Completar metadatos, fecha, protocolo y condiciones de uso de la
  fuente MH.
\item
  Obtener cartografía actualizada de drenaje y tratamiento del organismo
  operador.
\item
  Definir una campaña de verificación de establecimientos y descargas.
\item
  Acordar protocolo de aforo, muestreo, laboratorio y cadena de
  custodia.
\item
  Validar qué registros sectoriales corresponden a lavanderías
  industriales reales.
\item
  Definir criterios de evaluación económica, ambiental y de gobernanza.
\item
  Crear un entorno de software fijado para reproducir todo el
  repositorio en otro equipo.
\end{itemize}

\subsection{Registro resumido de hitos
Git}\label{registro-resumido-de-hitos-git}

{\def\LTcaptype{none} % do not increment counter
\begin{longtable}[]{@{}p{0.20\linewidth}p{0.44\linewidth}p{0.28\linewidth}@{}}
\toprule\noalign{}
Fecha & Commit & Hito \\
\midrule\noalign{}
\endhead
\bottomrule\noalign{}
\endlastfoot
25 mayo 2026 & \texttt{2eb2057} & Inicio del análisis geoespacial
reproducible. \\
25 mayo 2026 & \texttt{54d24f2} & Incorporación de auditoría. \\
11 junio 2026 & \texttt{5fab10d} & Avance de controles auditables. \\
25 junio 2026 & \texttt{6b640de} & Comparación con MH. \\
2 julio 2026 & \texttt{f284fa6} & Organización de productos legacy. \\
27 agosto 2026 & \texttt{e5831b0} & Pipeline textil e hídrico auditable
de Santa Ana. \\
24 septiembre 2026 & \texttt{3d69baf} & Reorganización de la fase
histórica en \texttt{version\_01}. \\
24 septiembre 2026 & \texttt{7cd11ab} & Integración de la fase moderna
regional. \\
24 septiembre 2026 & \texttt{a72c504} & Reporte metodológico y
siguientes pasos. \\
24 septiembre 2026 & \texttt{991d1ff} & Reporte sintético y ajustes
metodológicos. \\
24 septiembre 2026 & \texttt{e6c43c0} & README general del
repositorio. \\
\end{longtable}
}

\subsection{Nota para entrega
institucional}\label{nota-para-entrega-institucional}

Esta bitácora describe actividades y productos comprobables en el
repositorio. Para convertirla en un formato oficial de servicio social
deben validarse las 380 horas estimadas y agregarse, cuando la institución
lo requiera, las fechas exactas por jornada, firma o visto bueno,
objetivos del programa y relación de cada actividad con las competencias
o metas oficiales.

\end{document}
